\section*{الفصل \ref{ch:b3:palaeomagnetism} --- المغناطيسية الأرضية والمغناطيسية القديمة}

\begin{solution}{exo:b3:palaeomagnetism:1}
$\tan I = 2\tan\ang{30} = 1.1547$، إذن $I = \ang{49.1}$. وبالعكس، $\tan\lambda
= \tfrac12\tan(\ang{-20}) = -0.1820$، فإنّ $\lambda = \ang{-10.3}$، أي
\ang{10.3}\,جنوبًا.
\end{solution}

\begin{solution}{exo:b3:palaeomagnetism:2}
$H = \sqrt{20000^2 + 2000^2} = \qty{20100}{nT}$;
$F = \sqrt{20000^2 + 2000^2 + 40000^2} = \qty{44770}{nT}$. الانحراف
المغناطيسي هو $\arctan(-2000/20000) = \ang{-5.7}$، \ang{5.7} إلى الغرب من الشمال؛
\omterm{def:b3:palaeomagnetism:inclination}{الميل المغناطيسي} $\arctan(40000/20100) = \ang{63.3}$، نحو الأسفل.
\end{solution}

\begin{solution}{exo:b3:palaeomagnetism:3}
$B_0 = -g_1^0 = \qty{29350}{nT}$ عند خط الاستواء، و$2B_0 = \qty{58700}{nT}$ عند
القطبين. $m = 4\pi a^3B_0/\mu_0 = a^3B_0/\num{e-7} = \num{2.5862e20}\times
\num{2.935e-5}/\num{e-7} = \qty{7.59e22}{A.m^2}$، أقل قليلًا من
\qty{7.69e22}{A.m^2} لثنائي القطب الكامل، الذي يضيف $g_1^1$ و$h_1^1$.
\end{solution}

\begin{solution}{exo:b3:palaeomagnetism:4}
(أ) \omterm{def:b3:palaeomagnetism:nrm}{المغنطة المتبقية الحرارية}، تسجّل المجال أثناء برودة الحمم
عبر \omterm{def:b3:palaeomagnetism:blocking}{درجات حرارة الحجب} لحبيباتها. (ب) \omterm{def:b3:palaeomagnetism:nrm}{المغنطة المتبقية
الفتاتية}، تسجّل المجال خلال الزمن الذي استغرقه الغرين كي يستقرّ،
ويتثبّت بالانضغاط. (ج) \omterm{def:b3:palaeomagnetism:nrm}{المغنطة المتبقية الكيميائية}، تسجّل المجال
أثناء نموّ الهيماتيت، أي أثناء الدفن المبكّر، ومن ثمّ بعد الترسيب إلى حدٍّ ما.
(د) المغنطة المتبقية اللزجة وفق المجال الحالي، المضافة إلى
TRM الأصلية بفعل الحبيبات ذات أقصر زمن استرخاء؛
وتُزال بأولى خطوات إزالة المغنطة.
\end{solution}

\begin{solution}{exo:b3:palaeomagnetism:5}
متجهات الوحدة $(\cos I\cos D,\ \cos I\sin D,\ \sin I)$ هي
\begin{gather*}
  (0.6545, 0.1391, 0.7431),\quad (0.6670, 0.2695, 0.6947),\\
  (0.5235, 0.1501, 0.8387),\quad (0.5985, 0.0629, 0.7986).
\end{gather*}
المجموع هو $(2.4435, 0.6216, 3.0751)$، وطوله $R = 3.9766$. المتوسط:
$D = \arctan(0.6216/2.4435) = \ang{14.3}$, $I = \arcsin(3.0751/3.9766) =
\ang{50.7}$. $k = 3/(4 - 3.9766) = 128$. مع $20^{1/3} - 1 = 1.7144$،
$\cos\alpha_{95} = 1 - (0.0234/3.9766)\times1.7144 = 0.98991$، إذن
$\alpha_{95} = \ang{8.1}$: مع أربع عيّنات يبقى المخروط واسعًا، مهما
كان تماسك المجموعة.
\end{solution}

\begin{solution}{exo:b3:palaeomagnetism:6}
\sloppy $\lambda = \arctan(\tfrac12\tan\ang{50})$، أي $\ang{30.79}$، إذن $p = \ang{59.21}$.
ثم $\sin\lambda_p = \sin\ang{40}\cos\ang{59.21}$ زائد $\cos\ang{40}\sin\ang{59.21}
\cos\ang{20}$، أي $0.9474$، و$\lambda_p = \ang{71.34}$\,شمالًا. $\sin\beta =
\sin\ang{59.21}\sin\ang{20}/\cos\ang{71.34}$، أي $0.9184$، و$\beta = \ang{66.66}$.
لأنّ $\cos p = 0.5119$ أصغر من $\sin\ang{40}\times0.9474 = 0.6090$، تنطبق الحالة الثانية:
$\varphi_p = 10 + 180 - 66.66 = \ang{123.34}$\,شرقًا. يقع القطب عند
\ang{71.3}\,شمالًا، \ang{123.3}\,شرقًا.
\end{solution}

\begin{solution}{exo:b3:palaeomagnetism:7}
(أ) $(7.69 - 8.32)/8.32 = \qty{-7.6}{\percent}$ خلال 125 سنة، أو
\qty{-6.1}{\percent} لكلّ قرن. (ب) $29350/12.6 = 2329$، أي نحو 2300 سنة.
ليس هذا تنبّؤًا: إذ يتغيّر المعدّل نفسه من نموذج مدّته خمس سنوات إلى
الذي يليه، فقد كان العزم في الماضي أعلى وأدنى دون أن
ينعكس المجال على فترات غير منتظمة لا
يتنبّأ بها أيّ معدّل.
\end{solution}

\begin{solution}{exo:b3:palaeomagnetism:8}
بين 20 و\qty{300}{\celsius} فقد المتجه $(6.4, 0.0, 13.6)$: مركّبةً
بمقدار \qty{15.0}{mA/m}، مع $D = \ang{0.0}$
و$I = \arctan(13.6/6.4) = \ang{64.8}$. وفوق \qty{300}{\celsius} فقد $(25.8, 18.1, 24.6)$:
\qty{40.0}{mA/m}، $D = \arctan(18.1/25.8) = \ang{35.1}$، و$I =
\arctan(24.6/31.5) = \ang{38.0}$. أمّا الثانية، التي تمرّ بالأصل، فهي
ChRM؛ أمّا الأولى فتقع على اتّجاه شمالي شديد الانحدار، أُزيلت عند
درجة حرارة منخفضة، وهي تراكب لزج وفق المجال الحديث.
\end{solution}

\begin{solution}{exo:b3:palaeomagnetism:9}
$\lambda = \arctan(\tfrac12\tan(\ang{-35})) = \ang{-19.3}$. من أجل $I = \ang{-40}$
و\ang{-30}: \ang{-22.8} و\ang{-16.1}، أي فترة مقدارها \ang{6.7}.
$p = \ang{90} - \lambda = \ang{109.3}$, $\cos^2p = 0.1096$:
$\mathrm dp = 2.5\times(1 + 0.329) = \ang{3.3}$ و$\mathrm dm =
5\times\sin\ang{109.3}/\cos\ang{35} = \ang{5.8}$.
\end{solution}

\begin{solution}{exo:b3:palaeomagnetism:10}
(أ) $\tan\lambda = \tfrac12$، $\lambda = \ang{26.6}$. $F$ هو $B_0$ عند
خط الاستواء و$2B_0$ عند القطبين: نسبتها 2. (ب) $\tan I = 2\tan\lambda$
يعطي $(1 + \tan^2I)\,\mathrm dI = 2(1 + \tan^2\lambda)\,\mathrm d\lambda$، إذن
\[
  \frac{\mathrm dI}{\mathrm d\lambda} = \frac{2/\cos^2\lambda}{1 + 4\tan^2\lambda}
  = \frac{2}{\cos^2\lambda + 4\sin^2\lambda} = \frac{2}{1 + 3\sin^2\lambda}.
\]
خطأ $\delta I$ يعطي $\delta\lambda = \tfrac12(1 + 3\sin^2\lambda)\,\delta I$،
من $\tfrac12\delta I$ عند خط الاستواء إلى $2\,\delta I$ عند القطبين: أكبر
الخطأ يقع قرب القطبين. (ج) قرب خط الاستواء يتغيّر \omterm{def:b3:palaeomagnetism:inclination}{الميل المغناطيسي}
ضعف سرعة تغيّر خط العرض، فإنّ الخطأ المعطى في الميل يكلّف نصف
قدر من خط العرض؛ وقرب القطبين يكاد \omterm{def:b3:palaeomagnetism:inclination}{الميل المغناطيسي} لا يتغيّر مع خط العرض.
\end{solution}

\begin{solution}{exo:b3:palaeomagnetism:11}
(أ) $\cos p = \sin(\ang{-35})\sin\ang{70} + \cos\ang{35}\cos\ang{70}\cos\ang{150}
= -0.5390 - 0.2426 = -0.7816$، $p = \ang{141.41}$. ثمّ $\tan I_x = 2\cot p =
-2.506$، $I_x = \ang{-68.2}$، و$\lambda_x = \ang{90} - p = \ang{-51.4}$.
$\cos D_x = (\sin\ang{70} - \sin(\ang{-35})\cos p)/(\cos\ang{35}\sin p) =
0.9617$، ولأنّ القطب يقع شرق خط زوال الموقع ($\sin\ang{150} >
0$)، $D_x = \ang{15.9}$. (ب) $\lambda_o = \arctan(\tfrac12\tan(\ang{-60})) =
\ang{-40.9}$. $R = 30.0 - 15.9 = \ang{14.1}$: دوران الكتلة مع عقارب الساعة
حول محور رأسي. $P = \lambda_x - \lambda_o = -51.4 + 40.9 =
\ang{-10.5}$: كان التكوين يقع \ang{10.5} أبعد شمالًا مما تقتضيه الصفيحة المستقرّة،
ولذلك تحرّكت الكتلة منذ ذلك الحين بنحو \ang{10.5}،
أي نحو \qty{1170}{km}، جنوبًا بالنسبة إلى الصفيحة.
\end{solution}

\begin{solution}{exo:b3:palaeomagnetism:12}
(أ) بفرض $\mathrm d\Omega = 2\pi\sin\theta'\,\mathrm d\theta'$،
\[
  \int_0^\theta\frac{\kappa e^{\kappa\cos\theta'}}{4\pi\sinh\kappa}\,2\pi\sin\theta'\,
  \mathrm d\theta' = \frac{e^\kappa - e^{\kappa\cos\theta}}{2\sinh\kappa}
  = \frac{e^\kappa - e^{\kappa\cos\theta}}{e^\kappa - e^{-\kappa}} .
\]
(ب) من أجل $\kappa$ كبيرة، يكون $e^{-\kappa}$ مهملًا؛ وبالقسمة على $e^\kappa$
يبقى $1 - e^{-\kappa(1 - \cos\theta)} \approx 1 - e^{-\kappa\theta^2/2}$. وهذه النسبة
تساوي $1 - e^{-1} = \qty{63}{\percent}$ عندما $\kappa\theta^2/2 = 1$، $\theta =
\sqrt{2/\kappa}$ راديان، أو $\ang{81}/\sqrt\kappa$. (ج) $81/\sqrt{50} =
\ang{11.5}$؛ والصيغة الدقيقة تعطي \qty{63.1}{\percent}.
\end{solution}

\begin{solution}{pb:b3:palaeomagnetism:1}
\textbf{1.} يمكن أن ينمو الهيماتيت في الطبقات الحمراء من أكاسيد الحديد وهيدروكسيداته
بعد زمن طويل من ترسّب الرواسب، فتكتسب \omterm{def:b3:palaeomagnetism:nrm}{المغنطة المتبقية الكيميائية} من ذلك
العمر المتأخّر. قد يدلّ اختبار الطيّ الإيجابي، إذا طُويت الطبقات بعد وقت قصير
من الترسّب، أو اختبار الانعكاس الإيجابي عبر طبقات متقاربة من القطبيّتين معًا،
على أنّ المغنطة المتبقية مبكّرة.

\textbf{2.} $(\cos\ang{37}\cos\ang{50}, \cos\ang{37}\sin\ang{50}, \sin\ang{37})
= (0.5134, 0.6118, 0.6018)$.

\textbf{3.} $R = \sqrt{3.8943^2 + 2.6500^2 + 3.5717^2} = 5.9114$. $D =
\arctan(2.6500/3.8943) = \ang{34.2}$, $I = \arcsin(3.5717/5.9114) = \ang{37.2}$.

\textbf{4.} $k = 5/(6 - 5.9114) = 56.5$. ومع $20^{1/5} - 1 = 0.8206$،
$\cos\alpha_{95} = 1 - (0.0886/5.9114)\times0.8206 = 0.98771$, $\alpha_{95} =
\ang{9.0}$.

\textbf{5.} B: $k = 6/(7 - 6.8979) = 58.8$؛ $20^{1/6} - 1 = 0.6475$،
$\cos\alpha_{95} = 1 - (0.1021/6.8979)\times0.6475 = 0.99042$, $\alpha_{95} =
\ang{7.9}$. C: $k = 7/(8 - 7.8707) = 54.1$؛ $20^{1/7} - 1 = 0.5340$،
$\cos\alpha_{95} = 1 - (0.1293/7.8707)\times0.5340 = 0.99123$, $\alpha_{95} =
\ang{7.6}$.

\textbf{6.} الموقع B له أكبر $k$، أي أشدّ المجموعات تكتّلًا؛ والموقع C له أصغر
$\alpha_{95}$، أي أفضل المتوسطات تحديدًا، لأنّه يضمّ أكبر عدد من المواقع. وهذه الثلاثة
متقاربة: فالتكوينات الثلاثة جميعها مجمّعة جيدًا.

\textbf{7.} $\lambda_A = \arctan(\tfrac12\tan\ang{37.17}) = \ang{20.76}$\,شمالًا،
$\lambda_B = \ang{32.15}$\,شمالًا، $\lambda_C = \ang{45.33}$\,شمالًا.

\textbf{8.} $p = 90 - 20.76 = \ang{69.24}$. $\sin\lambda_p =
\sin\ang{47}\cos\ang{69.24} + \cos\ang{47}\sin\ang{69.24}\cos\ang{34.23} =
0.2593 + 0.5272 = 0.7865$: $\lambda_p = \ang{51.86}$\,شمالًا.

\textbf{9.} $\sin\beta = \sin\ang{69.24}\sin\ang{34.23}/\cos\ang{51.86} =
0.8516$، $\beta = \ang{58.40}$. لأنّ $\cos p = 0.3545 < \sin\ang{47}\times
0.7865 = 0.5752$، فتنطبق الحالة الثانية، $\cos\beta < 0$:
$\varphi_p = 15.00 + 180 - 58.40 = \ang{136.60}$\,شرقًا.

\textbf{10.} $\cos^2p = 0.1257$: $\mathrm dp = 4.50\times1.377 = \ang{6.2}$;
$\mathrm dm = 8.99\times\sin\ang{69.24}/\cos\ang{37.17} = \ang{10.5}$.

\textbf{11.} من \ang{51.9}\,شمالًا عند \qty{60}{Ma} إلى \ang{66.5}\,شمالًا عند
\qty{30}{Ma} و\ang{87.1}\,شمالًا عند \qty{4}{Ma}، على خط الزوال نفسه تقريبًا
(\ang{136}--\ang{139}\,شرقًا): مسار متّجه نحو القطب الجغرافي، يبلغه
ضمن مخروطه أحدث التكوين.

\textbf{12.} تقع طبقة الرماد عند \qty{22}{m} في النطاق الثاني ذي القطبية الطبيعية،
19.7--\qty{24.6}{m}؛ ومع عمر يقارب \qty{3.2}{Ma} يكون ذلك النطاق C2An.2n
(3.207--3.116). ثمّ إنّ النطاق 4.2--\qty{14.8}{m} هو C2An.3n، والنطاق
29.2--\qty{54.0}{m} هو C2An.1n، وتنتمي القاعدة المعكوسة إلى الزمن المغناطيسي الواقع أسفل
C2An، وتنتمي القمّة المعكوسة إلى الزمن المغناطيسي الواقع فوقه.

\textbf{13.} $(19.7 - 4.2)/(3.596 - 3.207) = 15.5/0.389 = \qty{39.8}{m/Myr}$;
$(54.0 - 24.6)/(3.116 - 2.581) = 29.4/0.535 = \qty{55.0}{m/Myr}$.

\textbf{14.} القاعدة: $3.596 + 4.2/39.8 = \qty{3.70}{Ma}$. والقمّة: $2.581 -
10.0/55.0 = \qty{2.40}{Ma}$.

\textbf{15.} لا يُؤرَّخ إلا حدود الانعكاس، التي تفصل بينها بضع مئات الآلاف من السنين؛
وبين كل حدّين يُفترض أنّ المعدّل ثابت. وسيغيّره تغيّر الهبوط،
أو تغيّر إمداد الرواسب أو تغيّر المناخ، أو فجوة في الترسيب، وقد تُزيل الفجوةُ الترسيبية
حتّى نطاقًا قصيرًا، فتُضلّل المطابقة.

\textbf{16.} $\cos p = \sin\ang{47}\sin\ang{78} + \cos\ang{47}\cos\ang{78}
\cos\ang{175} = 0.7154 - 0.1413 = 0.5741$, $p = \ang{54.96}$. $\tan I_x =
2\cot p = 1.4023$، $I_x = \ang{54.5}$؛ $\lambda_x = \ang{35.04}$\,شمالًا.

\textbf{17.} $\cos D_x = (\sin\ang{78} - \sin\ang{47}\times0.5741)/
(\cos\ang{47}\sin\ang{54.96}) = 0.99976$؛ يقع القطب شرق
خط زوال الموقع ($\sin\ang{175} > 0$)، ومنه $D_x = \ang{1.3}$.

\textbf{18.} خطوط العرض القديمة المتوقَّعة من $I_x$ = \ang{54.5}، \ang{60.0}
و\ang{63.6}: \ang{35.03}، \ang{40.89} و\ang{45.21}. $P_A = 35.03 - 20.76 =
\ang{14.3}$, $P_B = 40.89 - 32.15 = \ang{8.7}$, $P_C = 45.21 - 45.33 =
\ang{-0.1}$.

\textbf{19.} $R_A = 34.2 - 1.3 = \ang{32.9}$, $R_B = 23.9 - 2.9 = \ang{21.0}$,
$R_C = 3.4 - 2.5 = \ang{0.9}$، كلّها مع عقارب الساعة.

\textbf{20.} $\Delta D_o = \arcsin(\sin\ang{8.99}/\cos\ang{37.17}) = \ang{11.3}$;
$\Delta D_x = \arcsin(\sin\ang{3.0}/\cos\ang{54.5}) = \ang{5.2}$; $\Delta R_A =
\sqrt{11.3^2 + 5.2^2} = \ang{12.4}$.

\textbf{21.} $R_A = \ang{32.9} \pm \ang{12.4}$ و$R_B = \ang{21.0} \pm
\ang{13.8}$ دالّتان إحصائيًّا؛ أمّا $R_C = \ang{0.9} \pm \ang{17.9}$ فغير دالّ.

\textbf{22.} قبل ستّين مليون سنة كانت كيستا تقع بنحو \ang{14} من خطّ العرض،
جنوبَ الموقع الذي تضعه فيه القارة المستقرّة، وكانت قد أُدِيرت بنحو \ang{33}
عكس عقارب الساعة بالنسبة إلى توجّهها الحالي. وبحلول \qty{30}{Ma} كانت قد تحرّكت شمالًا بنحو \ang{6}
من أصل \ang{14}، ودارت بنحو \ang{12} من أصل \ang{33}
مع عقارب الساعة؛ وبحلول \qty{4}{Ma} بلغت موقعها الحالي على القارة،
دون أيّ دوران متبقٍّ. وانجرفت شمالًا على طول هامش القارة، وهي تدور
مع عقارب الساعة أثناء ذلك، واستقرّت في موضعها بحلول \qty{4}{Ma}.

\textbf{23.} $47.00 - 20.76 = \ang{26.24}$، أو $26.24\times111.19 =
\qty{2918}{km}$. يشمل ذلك حركة القارة المستقرّة نفسها؛
بالنسبة إليها، تحرّكت كيستا بمقدار $P_A = \ang{14.3}$.

\textbf{24.} $I = \ang{28.18}$ و\ang{46.16} تعطيان $\lambda = \ang{15.00}$
و\ang{27.50}: نصف عرض قدره \ang{6.25}، والانحراف المعياري الواحد
هو $6.25/1.96 = \ang{3.19}$، أو \qty{355}{km}. ومنذ التكوين A، انجرفت كيستا
شمالًا بمقدار
\textbf{\qty[text-series-to-math=true]{2918 \pm 355}{km}}.
\end{solution}
